\section{Writing Plans（实施计划）}

\subsection{目标与原则}

Writing Plans 解决的是 \textbf{How}——"如何实现"。它将宏观设计转化为\textbf{零上下文可执行}的细粒度步骤。

\begin{importantbox}{三大准则}
\begin{enumerate}
    \item \textbf{DRY}（Don't Repeat Yourself）：避免写重复代码和重复逻辑。
    \item \textbf{YAGNI}（You Aren't Gonna Need It）：不做"可能以后用得上"的功能。
    \item \textbf{TDD}：先写测试再写实现，频繁提交、小步迭代。
\end{enumerate}
\end{importantbox}

\subsection{产出物：Plan.md}

计划文档（\texttt{Plan.md}）包含极其具体的细节：
\begin{itemize}
    \item 要创建哪个确切路径的文件。
    \item 要修改哪几行。
    \item 要写什么样的失败测试。
    \item 要运行什么确切的命令来验证。
\end{itemize}

\begin{knowledgebox}{消除上下文依赖}
计划的编写标准是：假设最终写代码的工程师（或另一个 Subagent）对项目一无所知，仅凭此计划即可完成开发。因此可以另开一个会话，让 Codex / Claude Code / 人类工程师直接执行。
\end{knowledgebox}

\subsection{计划完成后的选择}

计划写完后不写代码，而是询问用户：
\begin{itemize}
    \item 在当前会话中渐进式执行（Sequential Development）。
    \item 开新会话批量执行（Parallel Agents）。
\end{itemize}

\subsection{本章小结}

Writing Plans 将架构设计"翻译"为任何人/Agent 都能执行的微观步骤清单，是从"想"到"做"的关键桥梁。


